% A longer document, in chapters.
%
% The difference from the article template is structure rather than
% packages: a report has chapters, a table of contents that is worth having,
% and each chapter in its own file so that a build of one chapter is cheap.
% NextTex builds the chapter you are typing in and not the whole document,
% which is what makes a hundred-page report editable at all.

\documentclass[11pt]{report}

\usepackage[T1]{fontenc}
\usepackage{lmodern}
\usepackage[margin=1in]{geometry}
\usepackage{parskip}
\usepackage{microtype}

\usepackage{amsmath, amssymb}
\usepackage{graphicx}
\usepackage{booktabs}
\usepackage{siunitx}

\usepackage[backend=biber, style=numeric, sorting=none]{biblatex}
\addbibresource{references.bib}

\usepackage{hyperref}
\hypersetup{colorlinks, linkcolor=black, citecolor=black, urlcolor=blue}

\title{A Working Title}
\author{Your Name}
\date{\today}

\begin{document}
\maketitle
\tableofcontents

% One file per chapter. Click a chapter in the file list and type in it:
% NextTex works out which document owns it and builds that, so the whole
% report is not retypeset on every keystroke.
\chapter{Introduction}
\label{ch:introduction}

Start here. A citation looks like this~\cite{knuth1984}, and a reference to
another chapter looks like Chapter~\ref{ch:method}.

\section{Background}

What somebody needs to know before the rest of this makes sense.

\chapter{Method}
\label{ch:method}

What was done, in enough detail that it could be done again.

\begin{equation}
  \label{eq:method}
  E = mc^{2}.
\end{equation}

Equation~\eqref{eq:method} is referred to from Chapter~\ref{ch:introduction}
by its label rather than by its number, so inserting a chapter above it
changes nothing here.


\printbibliography

\end{document}
